\documentclass[11pt, a4paper]{article}

\usepackage[utf8]{inputenc}
\usepackage{amsmath, amssymb}
\DeclareMathOperator{\diag}{diag} 
\usepackage{graphicx}
\usepackage{hyperref} 

\title{A connectivity modelling toolset for ecological and geospatial edge computing in WebAssembly/rust}
\author{Alex Southgate}
\date{\today}

\begin{document}

\maketitle

\begin{abstract}
Circuit theory has been successfully applied to ecological connectivity modelling, notably via the Circuitscape software, which is typically run locally on a laptop or via a server. For downstream geospatial web applications relying on connectivity analysis, backend infrastructure is required, which can be costly. Recent developments in WebAssembly now allow fast C++ or Rust code to be run directly in a sandboxed browser environment for edge computing. We present a WebAssembly/Rust toolset with a geospatial data pipeline and efficient edge-computing implementation of connectivity analysis. This approach may be useful for geospatial modelling software where rasters and memory footprint are small enough for the browser context. Our implementation offers warm starts for repeated runs in interactive analysis, such as for repeated resistance map editing (e.g. adding features in scenario planning). Our results show that as expected, run-times for 1000x1000 raster networks are slower than Circuitscape, but fast enough to run in the browser without requiring a server.
\end{abstract}

\section{Introduction}

The theory of linear circuits has found useful application in ecological connectivity modelling, 
population genetics, and conservation \cite{mcrae2008using, mcrae2006isolation, mcrae2007circuit}, 
especially through the Circuitscape \cite{mcrae2016circuitscape} software. 
Circuit theory offers advantages beyond methods such as least-cost paths (see \cite{mcrae2007circuit}).
At the same time, these modelling approaches can require large compute resources, cloud instances or servers,
and technical expertise for maintenance.

In modern web applications, 
a front-end browser application is often paired with a backend server which performs more computationally expensive jobs. 
When scientific code is deployed into production web applications for practical use by customers (e.g. analysts) 
outside of academia, a backend is required. However, for many applications, maintaining this infrastructure is costly. 
Edge computing means making use of compute close to the `edge' of the cloud, 
which is either close to the user's computer, or the user's computer itself. 
The problem with performing complex calculations in the browser is that browser code 
(i.e. JavaScript) execution is strictly controlled for security purposes, which has poor performance and filesystem interaction.
To resolve this problem, WebAssembly was recently introduced,
allowing fast C++ or Rust code to be run in a sandboxed browser environment. 
It has previously been argued that WebAssembly is poised to transform scientific computing \cite{perkel2024no}.
In addition, web applications may integrate complex user controls
(e.g. modelling controls directly on a map) that are time consuming to complete using scientific scripts,
requiring expertise not always available to users outside of academia.

WebAssembly has found use in several scientific projects, including: Pyodide, a CPython port, supporting installs of NumPy and SciPy, amongst other packages\cite{pyodide2026};
WebR \cite{webr2026}, a WebAssembly compilation of the R language; ITK-Wasm, for spatial/image analysis \cite{itkwasm2026}. 
Application domains include software for mapping pipeline corrosion \cite{guarneri2024web}, viral genomics \cite{ji2024viralwasm}, or cancer genome visualisation \cite{finney2018chromatic}.


This paper introduces a Rust library and WebAssembly bindings to allow connectivity analysis directly within web applications,
such as geospatial analysis software. We demonstrate reasonable
 performance provided input images fit within browser sandbox memory limits,
 which has potential use for deployed scientific applications.

\section{Methodology}

\subsection{Vector rasterisation}

In order to keep the WASM toolset lightweight, we avoid heavy geospatial dependencies. For vector rasterisation, we use
the \texttt{geojson} and \texttt{geo} crates, which are lightweight and have no external dependencies. Handling vector inputs allows
applications with vector data to use the WASM toolset directly, and interface cleanly with vector tile
data such as MVT derived from mbtiles.

\subsection{Example datasets}

In order to demonstrate the functionality of our WASM toolset, a test region corresponding to a 1~$\text{km}^2$ region centered at
286600E, 79050N (EPSG:27700) was selected. The region includes buildings, roads, and a river. Vector data was obtained from
the UK Ordnance Survey open data \cite{os_data}. 
DSM and DTM LIDAR data was obtained from the UK Environment Agency \cite{ea_lidar_composite_2022}.

\subsection{Benchmarking}

Benchmarking was undertaken on a laptop with an AMD Ryzen 7 5700U, 16G of RAM, Linux 6.8.0-51-generic, and using Google Chrome 150.0.7871.46.
Benchmarking code was wrapped in a React applet which can be run locally and accessed via the code repository. Two benchmarks were run. 
In the first case, a resistance map was generated at 1000$\times$1000 resolution, then downsampled to lower resolutions, and the analysis code run with 5 repeats each.
It was expected that, in the browser context, the first run would be slowest due to V8/WASM JIT compilation engine warm-up. This is the first type of `cold-start'
distinct from the `warm-start' algorithm approach by reusing voltage rasters. In the second case,
the 1000x1000 raster was repeatedly edited by the addition of a randomly placed building; the time taken for post-edit runs (with `warm-start' initialization) was 
compared to runs executed without the `warm-start' (i.e. $v=0$ initialization) algorithm approach to evaluate the ability to reuse solutions for repeated runs.

\subsection{Connectivity model}

We use the methodology introduced by \cite{mcrae2006isolation} to describe connectivity with circuit theory, and a brief summary follows. First, Ohm's law \cite[Ch. 2]{alexander2007fundamentals} is \eqref{eq:ohm}:

\begin{equation}
\label{eq:ohm}
I = CV
\end{equation}

Secondly, for a given node (pixel), the conservation equation (Kirchhoff's equation, see \cite[Ch. 2]{alexander2007fundamentals}) \eqref{eq:conservation} just means that movement in is equal to movement out. This does treat the modelled species as a conserved quantity and the simplest approach ignores e.g. death. Taking $S_i$ as a source at node $i$:

\begin{equation}
\label{eq:conservation}
I_i = \sum_{j \in N(i)} C_{ij}V_{ij} - S_i = 0
\end{equation}

Given these two equations, only voltage differences $V_{ij}$ are unknown. Voltage, as a potential can be assumed to be a scalar quantity for a single node up to some constant, which leads to a system of equations \eqref{eq:system}:

\begin{equation}
\sum_{j \in N(i)} C_{ij}(V_i - V_j) = S_i
\end{equation}

\begin{equation}
\Big(\sum_{j \in N(i)} C_{ij}\Big)V_i - \sum_{j \in N(i)} C_{ij}V_j = S_i
\end{equation}

\begin{equation}
\label{eq:system}
\rightarrow \Big(\mathbf{D_C}- \mathbf{M}\Big)\mathbf{v} = \mathbf{L}\mathbf{v} = \mathbf{s}
\end{equation}

Which is a discrete Poisson problem \cite[Ch. 4]{leveque2007finite}. The matrix $\mathbf{L}$ is the Laplacian. The diagonal elements of $\mathbf{L}$ are precomputed. For ground (sink) nodes, either the corresponding rows can be ignored in calculation and set to zero, or the diagonal elements set to $1$ and off-diagonal to zero. $\mathbf{L}$ is sparse, since only neighboring nodes are connected in the simplest case.

\subsection{Modes}

The simplest run mode is for a source and sink pair, equivalent to a conditional random walk (paths of those that travel from source to sink). If a raster image of source nodes is provided, then the resulting output is connectivity given inputs from the combined input sources, solved together, as in a conventional physical system like electrical circuits. Instead of the combined mode, where current will flow according to the global pressure gradient, some users may choose to intentionally superimpose $n \choose 2$ pairwise simulations from pairs of $n$ nodes, or $n \times m$ simulations between $n$ sources and $m$ sinks. Each approach has different assumptions and interpretation. Each mode has the same underlying mathematical system and solver; in the case of a single pair, only two rows of the matrix $L$ are source/sink; in the case of a raster of currents, many rows are set to sources. Superposition of the different cases can be undertaken in a loop over the chosen sets of pairs.

\subsection{Solver}

Given an $N\times M$ image, the matrix $L$ takes $O(NM)$ space. In Circuitscape, conjugate gradient with algebraic multigrid (CG+AMG) or a Cholesky solver is used via optimised CHOLMOD libraries that are not available in a WASM context. We utilise the conjugate gradient method with jacobi preconditioning for simplicity. This method is iterative and memory efficient, at $O(MN)$ space and $O(kMN)$ time for $k$ iterations. We implement our CG solver in order to reduce dependencies for the WASM context. 

A good description of the CG method can be found in \cite{shewchuk1994introduction} or \cite[Ch. 4]{leveque2007finite}. Since we implement it manually for WASM, a brief explanation follows. Solving the system:

\[
\mathbf{L}\mathbf{v} = \mathbf{s}
\]

is equivalent to solving the energy function:

\[
f(\mathbf{v}) = \frac{1}{2} \mathbf{v}^T \mathbf{L}\mathbf{v} - \mathbf{v}^T\mathbf{s}
\]

Since when taking the negative gradient and setting to zero, we get the residual error for our system:

\[
-\nabla f(\mathbf{v}) = \mathbf{s} - \mathbf{L}\mathbf{v} = \mathbf{r}
\]

Stepping in the direction $r$ at step $k$ gives us a new guess.

\[
\mathbf{v}_{k+1} = \mathbf{v}_{k} + \mathbf{\alpha} \mathbf{r}_k
\]

Taking the gradient and choosing an optimal step size $\alpha$ gets to the bottom of the hill for that slice. The next step will be orthogonal to the previous one, since if it were not, the previous one would not have been the furthest (best) possible step. In steepest descent, to get the best step, applying calculus to the energy function it can be shown (see \cite[Ch. 4]{leveque2007finite}) that:

\[
\mathbf{\alpha}_{k} = \frac{\mathbf{r}_{k}^T \mathbf{r}_k}{\mathbf{r}_k^T \mathbf{L}\mathbf{r}_k}
\]

However, stepping maximally in a sequence of orthogonal steps can be inefficient and lead to zig-zag steps. Instead, a kind of moving average or momentum can be applied to the steps. 

\[
\mathbf{p}_{k+1} = \mathbf{r}_{k+1} + \mathbf{\beta}_{k+1}\mathbf{p}_k
\]

In the case of CG, $\beta$ is:

\[
\beta_{k+1} = \frac{\mathbf{r}^T_{k+1}\mathbf{r}_{k+1}}{\mathbf{r}_{k}^T\mathbf{r}_{k}}
\]

We also implement a preconditioner. As described in \cite[Ch. 4.3.4]{leveque2007finite}, convergence rate depends on condition number. We solve a transformed system:

\[
\mathbf{M}^{-1}\mathbf{L}\mathbf{v} = \mathbf{M}^{-1}\mathbf{s}
\] 

Where $\mathbf{M}^{-1}$ approximates $\mathbf{L}$, 
reducing the condition number. 
We utilise a simple diagonal preconditioner with $\mathbf{M} = \diag{\mathbf{L}}$, 
also known as a Jacobi preconditioner, 
which can be useful in instances of highly variable parameters such as material density or irregular meshes \cite{wathen2015preconditioning}.
We utilise this simple method to reduce WASM dependencies and binary footprint. 
For front-end web applications, it is crucial that data transfer is minimised
For some ecological rasters, this method is not preferred,
such as in Circuitscape, where the algebraic multigrid (AMG) preconditioner is offered \cite{hall2021circuitscape}.

\subsection{Code}

Code can be obtained at \url{https://github.com/asouthgate/wasm-connectivity}. This code contains several simple examples. All results can be easily reproduced by running the React development server, clicking to the benchmark panel, running the benchmarks, and downloading the results.

\section{Results}

Figure \ref{fig:wasm-connectivity} gives an example of a 1000x1000 resistance input and current map output
rendered in a React frontend using raster source mode (`advanced mode') for the test region.
Figure \ref{fig:benchmarking} shows time and memory benchmarks as a function of raster size. 

\begin{figure}[h]
\centering
\includegraphics[width=0.8\textwidth]{figures/resistance_current_react.png}
\caption{Example dataset 1000x1000 resistance and current maps rendered using a React frontend with WASM connectivity library in 6 minutes.}
\label{fig:wasm-connectivity}
\end{figure}


\begin{figure}[h]
\centering
\includegraphics[width=0.8\textwidth]{figures/benchmarking.png}
\caption{Benchmarking results. Point in grey is cold start run, expected to be slower.}
\label{fig:benchmarking}
\end{figure}

\section{Discussion}

For front-end web applications, it is crucial that data transfer is minimised

We present a WebAssembly/Rust connectivity library for interactive geospatial edge-computing. This method can be integrated into geospatial web applications. Future work could be used to explore alternative solver methods and optimisations for the browser context. 

\section{LLM disclaimer}

This paper was written entirely by hand without LLM assistance. Code was developed with AI assistance. 

\bibliographystyle{plain} 
\bibliography{refs} 

\end{document}
